% Figura 4-3: Detalle hardware pasarela de borde RPi4 + módulo AHPI6108E HaLow
% Versión refactorizada: RPi4 con HaLow nativo vía SDIO (no Ethernet a router externo)

\begin{figure}[H]
\centering
\begin{tikzpicture}[
    box/.style={rectangle, draw=black!70, thick, rounded corners=3pt,
                minimum width=4.2cm, minimum height=2.2cm, align=center},
    chip/.style={rectangle, draw=#1!60, fill=#1!8, thick, rounded corners=2pt,
                 minimum width=2.8cm, minimum height=0.7cm, font=\footnotesize, align=center},
    chip/.default=blue,
    port/.style={rectangle, draw=black!50, fill=gray!20, thick,
                 minimum width=1.8cm, minimum height=0.35cm, font=\tiny, align=center},
    lbl/.style={font=\scriptsize, align=center},
    arr/.style={->, thick, >=stealth, draw=black!70},
    darr/.style={<->, thick, >=stealth, draw=black!70},
]

% ============ MÓDULO AHPI6108E (HAT) ============
\node[box, fill=cyan!6, minimum height=3.2cm, minimum width=5.2cm] (hat) at (0,2.8) {};
\node[above=2pt of hat.north, font=\small\bfseries, text=cyan!40!black] {Alfa Networks AHPI6108E (HAT)};

% Chip MM6108 dentro del HAT
\node[chip=cyan, minimum width=3.4cm] at (0,3.6) (mm) {\textbf{MM6108} \tiny --- 802.11ah HaLow};

% Socket M.2 E-Key
\node[lbl, font=\tiny, text=gray!60!black] at (0,2.8) {M.2/NGFF 2230 E-Key};

% Conector GPIO 40-pin (salida inferior del HAT)
\node[port, minimum width=3.4cm, fill=yellow!15] at (0,1.85) (gpio_hat) {GPIO 40-pin $\rightarrow$ SDIO};

% Antena
\node[lbl, text=red!50!black] at (4.2,3.6) (ant) {Antena SMA\\915\,MHz\\IPEX4/MHF4};
\draw[arr, red!50] (mm.east) -- ++(0.8,0) -- (ant.west);

% MikroBUS (futuro)
\node[lbl, text=gray!50!black, font=\tiny] at (-4.0,3.4) {mikroBUS\\(SPI, I\textsuperscript{2}C, UART)};
\draw[arr, gray!40, dashed] (mm.west) -- ++(-0.8,0);

% Specs label (nota lateral para no solapar conexión SDIO)
\node[draw=gray!30, fill=gray!4, rounded corners=2pt, font=\tiny,
  text=gray!60!black, inner sep=4pt, text width=3.4cm, anchor=north west] at (3.6, 1.5) {
    \textbf{Specs AHPI6108E:}\\
    65$\times$56$\times$19\,mm\\
    TX: +25\,dBm (MCS0)\\
    RX: $-$100\,dBm\\
    32,5\,Mbps PHY m\'ax.
};

% ============ RASPBERRY PI 4 ============
\node[box, fill=orange!4, minimum height=5.0cm, minimum width=5.2cm] (rpi) at (0,-4.0) {};
\node[above=2pt of rpi.north, font=\small\bfseries, text=orange!50!black] {Raspberry Pi 4 (Pasarela de borde)};

% GPIO receptor
\node[port, minimum width=3.4cm, fill=yellow!15] at (0,-1.9) (gpio_rpi) {GPIO 40-pin $\leftarrow$ SDIO};

% Chips internos
\node[chip=orange, minimum width=3.4cm] at (0,-2.9) (cpu) {\textbf{BCM2711} \tiny --- Cortex-A72 4$\times$1,5\,GHz};
\node[chip=blue, minimum width=3.4cm] at (0,-3.9) (docker) {\textbf{Docker} \tiny --- TB Edge + PostgreSQL + Redis};
\node[chip=teal, minimum width=2.2cm] at (0,-4.9) (nrf) {\textbf{nRF52840}};
\node[lbl, text=teal!60!black, font=\tiny] at (0,-5.4) {USB Dongle --- Thread OTBR (RCP)};

% Puerto Ethernet
\node[port, minimum width=2.6cm] at (0,-6.2) (eth_r) {ETH GbE (gestión)};

% ============ CONEXIÓN SDIO (HAT → RPi4) — en fondo ============
\begin{scope}[on background layer]
\draw[darr, cyan!60, line width=2pt] (gpio_hat.south) -- (gpio_rpi.north)
    node[midway, right=12pt, font=\footnotesize, text=cyan!50!black, align=left,
         fill=white, inner sep=4pt, rounded corners=1pt] {
        \textbf{Bus SDIO}\\[-1pt]
        \tiny Nativo, sin\\[-1pt]
        \tiny intermediarios
    };
\end{scope}

% ============ SEÑALES INALÁMBRICAS ============
% Thread
\node[font=\scriptsize, text=teal!60!black, align=center] at (4.6,-4.9) {Red Thread\\[-1pt]\tiny 802.15.4 @ 2,4\,GHz\\[-1pt]\tiny $\leq$15 nodos/pasarela};
\draw[arr, teal!60, dashed] (nrf.east) -- ++(1.8,0);

% HaLow mesh (desde la antena)
\node[font=\scriptsize, text=cyan!50!black, align=center] at (4.6,2.8) {Malla HaLow\\[-1pt]\tiny 802.11s + HWMP\\[-1pt]\tiny 902--928\,MHz};
\draw[arr, cyan!60, dashed] (ant.east) -- ++(0.6,0) -- ++(0,-0.7);

% ============ NOTA TRABAJO FUTURO ============
\node[rectangle, draw=gray!40, fill=gray!5, rounded corners=2pt, 
      font=\tiny, text=gray!60!black, text width=3.0cm, align=center] 
      at (-4.8,-4.5) {
    \textbf{Trabajo futuro:}\\
    Interfaz SPI alternativa\\
    Chipset MM8108 (USB)
};

\end{tikzpicture}
\caption{Detalle de \textit{hardware} de la pasarela de borde con HaLow nativo.  El módulo Alfa Networks AHPI6108E (formato HAT, 65$\times$56\,mm) se acopla al conector GPIO de 40~pines de la Raspberry~Pi~4, conectando el chipset Morse Micro MM6108 mediante interfaz \textbf{SDIO} nativa al SoC BCM2711.  Esta integración elimina la necesidad de un enrutador HaLow externo: la pasarela participa directamente en la malla 802.11s como \emph{Mesh Point}.  El nRF52840 USB Dongle proporciona la radio IEEE~802.15.4 para Thread (OTBR).  El \emph{socket} mikroBUS ofrece interfaces SPI, I\textsuperscript{2}C y UART para futuras evaluaciones de buses alternativos y del chipset MM8108.}
\label{fig:halow-hardware-detail}
\end{figure}
